%%%%%%%%%%%%%%%%%%%%%%%%%%%%%%%%%%%%%%%%%%%%%%%%%%%%%%%%%%%%%%%%%%%%%%
% Lớp 10 | Phương pháp tọa độ trong mặt phẳng (𝑂𝑥𝑦) | Tọa độ của véctơ
% Dạng 0H9?1-1: Tọa độ điểm, độ dài đại số của véctơ trên 1 trục
% Số câu: 13  (sinh tự động bởi tach_cau.py; từ nay thầy sửa trực tiếp file này)
%%%%%%%%%%%%%%%%%%%%%%%%%%%%%%%%%%%%%%%%%%%%%%%%%%%%%%%%%%%%%%%%%%%%%%

%%%%%%%%%%%%%%%% TRẮC NGHIỆM 4 PHƯƠNG ÁN %%%%%%%%%%%%%%%%

% nguồn: MaHoaNH_TracNghiem_Lop10
\begin{ex}%[0H9N1-1]%[ma:c034413a]
	Trong mặt phẳng $Oxy$, cho ba điểm $A(2;0)$, $B(-1;-2)$, $C(5;-7)$. Tọa độ trọng tâm của $\Delta ABC$  là
	 \choice
		 {$(2;3)$}
		{\True  $(2;-3)$}
		{$(3;2)$} 
		 {$(-3;2)$}
	\loigiai{ 
	Gọi $G$ là trọng tâm của $\Delta ABC$, ta có $\heva{&x_G=\dfrac{x_A+x_B+x_C}{3}=2 \\&y_G=\dfrac{y_A+y_B+y_C}{3}=-3}\Rightarrow  G(2;-3)$.}
\end{ex}

% nguồn: MaHoaNH_TracNghiem_Lop10
\begin{ex}%[0H9N1-1]%[ma:3787f24e]
	Trong mặt phẳng với hệ tọa độ $Oxy$, cho hai điểm $A(2;-5)$ và $B(4;1)$. Tọa độ trung điểm $I$ của đoạn thẳng  $AB$ là
	\choice 
		 {\True $I(3;-2)$}
		 {$I(3;2)$}
		{$I(1;3)$} 
		{$I(-1;-3)$} 
	 \loigiai{
	Tọa độ trung điểm $I$ của đoạn thẳng $AB$ là  $\heva{&x_I=\dfrac{x_A+x_B}{2} \\&y_I=\dfrac{y_A+y_B}{2}}\Leftrightarrow \heva{&x_I=3 \\&y_I=-2}\Rightarrow I(3;-2)$.}
\end{ex}

% nguồn: MaHoaNH_TracNghiem_Lop10
\begin{ex}%[0H9N1-1]%[ma:5837cec2]
	Trong mặt phẳng tọa độ $Oxy$,  cho hai điểm $A(1;4)$ và $B(3;5)$. Khi đó
	\choice 
		 {$\overrightarrow{AB}=(-2;-1)$}
		{\True $\overrightarrow{AB}=(2;1)$} 
		{$\overrightarrow{AB}=(4;9)$} 
		 {$\overrightarrow{AB}=(1;2)$}
	\loigiai{ 
		 $\overrightarrow{AB}=(3-1;5-4)=(2;1)$.
	} 
\end{ex}

% nguồn: MaHoaNH_TracNghiem_Lop10
\begin{ex}%[0H9H1-1]%[ma:00f81650]
	Trong hệ trục $\left(O;\overrightarrow{i};\overrightarrow{j}\right)$, cho $\overrightarrow{OM}=3\left(\overrightarrow{i}+4\overrightarrow{j}\right)+5\overrightarrow{j}-2\overrightarrow{i}$. Tọa độ  điểm $M$ là
	\choice 
		{$(-1;-9)$} 
		 {\True $(1;17)$}
		 {$(-1;-17)$}
		 {$(1;9)$}
	\loigiai{ 
		Ta có $\overrightarrow{OM}=3\left(\overrightarrow{i}+4\overrightarrow{j}\right)+5\overrightarrow{j}-2\overrightarrow{i}=\overrightarrow{i}+17\overrightarrow{j}$.  Suy ra $M(1;17)$.
	 }
\end{ex}

% nguồn: MaHoaNH_TracNghiem_Lop10
\begin{ex}%[0H9H1-1]%[ma:8bbfc44b]
	Trong mặt phẳng $Oxy$, cho các  điểm $A(0;1)$, $B(2;-1)$, $C(1;6)$ và $D(1;7)$. Tam giác nào sau đây có cùng trọng tâm với tam giác $ABC$?
	 \choice
		{\True $\triangle OBD$} 
		{$\triangle  ABD$}
		{$\triangle  BCD$}
		 {$\triangle OAB$}
	 \loigiai{
		Gọi $G$ là trọng tâm của tam giác $ABC$.\\ 
		Ta có $\heva{&x_G = \dfrac{x_A+x_B+x_C}{3}=\dfrac{0+2+1}{3}=1 \\&y_G = \dfrac{y_A+y_B+y_C}{3}=\dfrac{1+(-1)+6}{3}=2}$ $\Rightarrow G(1;2)$.\\ 
		Gọi $K$ là trọng tâm của  tam giác $OBD$.\\
		Ta có  $\heva{&x_K = \dfrac{x_O+x_B+x_D}{3}=\dfrac{0+2+1}{3}=1 \\&y_K = \dfrac{y_O+y_B+y_D}{3}=\dfrac{0+(-1)+7}{3}=2}$ $\Rightarrow K(1;2)$.\\
		 Vậy tam giác $OBD$ có cùng trọng tâm với tam giác $ABC$.
	 }
\end{ex}

% nguồn: MaHoaNH_TracNghiem_Lop10
\begin{ex}%[0H9H1-1]%[ma:b0cb02fc]
	Trong mặt phẳng $Oxy$, cho tam giác  $OAB$ có $O(0;0)$, $A(6;0)$, $B(0;8)$. Gọi $M$, $N$, $P$ lần lượt là trung điểm của các cạnh $OA$, $OB$, $AB$. Đường tròn đi qua ba điểm $M$, $N$, $P$ có bán kính là
	 \choice
		{\True  $\dfrac{5}{2}$}
		 {$5$}
		{$5\sqrt{2}$} 
		{$10$} 
	\loigiai{ 
	 \immini{
		\begin{itemize}
			\item Ta có $\triangle MPN$ vuông tại $P$ nên đường tròn ngọai tiếp $\triangle MPN$  có tâm là trung điểm $MN$ và bán kính $R=\dfrac{MN}{2}=\dfrac{AB}{4}$.
			\item Vì $AB=\sqrt{(-6)^2+8^2}=10$ nên  $R=\dfrac{10}{4}=\dfrac{5}{2}$.
		\end{itemize}
		 }{
		\begin{tikzpicture}[>=stealth,line join=round,line cap=round,font= \footnotesize,scale=0.75]
				 \path
					(0,0)  coordinate (O)
					(2,0) coordinate  (A)
					(0,4)  coordinate (B)
					($(O)!1/2!(A)$)  coordinate (M)
					($(O)!1/2!(B)$) coordinate (N) 
					 ($(A)!1/2!(B)$) coordinate (P);
				\draw (A)--(B)--(O)--cycle  (M)--(N)--(P)--cycle;
				\foreach \x/\g in {A/0,B/60,O/180,M/-90,N/180,P/0} 
				 \fill (\x) circle (1pt)
					($(\x)+(\g:3mm)$)  node{$\x$};
			\end{tikzpicture}
		 }
	} 
\end{ex}


%%%%%%%%%%%%%%%% ĐÚNG/SAI %%%%%%%%%%%%%%%%

% nguồn: MaHoaNH_DungSai_Lop10
\begin{ex}%[0H9N1-1]%[ma:edc73d8b]
	Trên hệ trục $(O,\vec{e})$, cho các điểm $A(-3)$, $B(2)$, $C(7)$. Xét tính đúng sai của  các mệnh đề sau
	\choiceTF 
		 {$\overline{AB}=-5$}
		{\True  $\overline{OA}=-3$}
		{$\overline{CB}=5$} 
		{\True $\overline{CO}=-7$} 
	\loigiai{ 
		\begin{itemchoice}
			\itemch Ta  có $\overline{AB}=2-(-3)=5$.
			\itemch  Ta có $\overline{OA}=-3-(0)=-3$.
			\itemch Ta  có $\overline{CB}=2-7=-5$.
			 \itemch Ta có $\overline{CO}=0-7=-7$.
		\end{itemchoice}
	} 
\end{ex}

% nguồn: MaHoaNH_DungSai_Lop10
\begin{ex}%[0H9N1-1]%[ma:42bfa119]
	Trên hệ trục $(O,\vec{e})$, cho các điểm $A(-4)$, $B(3)$, $C(5)$. Xét tính đúng sai của các mệnh  đề sau
	\choiceTF 
		{Trung điểm  của $AB$ là $O$.}
		{\True  $AB=7$}
		{\True Trung điểm của  $BC$ là $D(4)$}
		{\True $\overline{AC}=9$} 
	\loigiai{ 
		\begin{itemchoice}
			\itemch Trung  điểm của $AB$ là $E\left(-\dfrac{1}{2} \right) $.
			\itemch Ta có $AB=7$. 
			\itemch Ta có trung điểm của $BC$  là $D(4)$.
			 \itemch Ta có $\overline{AC}=9$.
		\end{itemchoice}
	} 
\end{ex}

% nguồn: MaHoaNH_DungSai_Lop10
\begin{ex}%[0H9N1-1]%[0H9H1-1]%[ma:0c4ed6d2]
	Trong mặt phẳng $Oxy$, cho ba điểm $A(3;5)$, $B(-1;0)$, $C(4;-2)$. Xét tính  đúng, sai của các mệnh đề sau
	\choiceTF 
		{\True Độ dài véc-tơ $\overrightarrow{AB}$ bằng  $\sqrt{41}$}
		 {$\overrightarrow{BC}=(5;-2)$}
		{$\overrightarrow{AB}\cdot \cdot\overrightarrow{AC}=-31$} 
		{\True Nếu điểm $M$ có  tọa độ là $(2;1)$ thì $\overrightarrow{MA}+\overrightarrow{MB}+\overrightarrow{MC}=\overrightarrow0$}
	 \loigiai{
		Gọi $(x;y)$ là tọa độ của điểm $M$.\\ 
		Ta có $\overrightarrow{MA}=(3-x;5-y)$,  $\overrightarrow{MB}=(-1-x;-y)$, $\overrightarrow{MC}=(4-x;-2-y)$.\\
		Vì $\overrightarrow{MA}+\overrightarrow{MB}+\overrightarrow{MC}=\overrightarrow0$ nên 
		$\heva{&3-x-1-x+4-x=0\\&5-y-y-2-y=0}\Leftrightarrow \heva{&x=2\\&y=1.}$\\ 
		 Ta có $\overrightarrow{BC}=(5;-2)$, $\overrightarrow{AB}\cdot\overrightarrow{AC}=(-4;-5)\cdot(1;-7)=31$.
		\begin{itemchoice}
			\itemch Độ dài véc-tơ $\overrightarrow{AB}$ bằng $\sqrt{41}$ đúng. 
			 \itemch $\overrightarrow{BC}=(5;-2)$ sai.
			\itemch Tích vô hướng của hai véc-tơ $\overrightarrow{AB}$ và $\cdot\overrightarrow{AC}=-31$ sai. 
			 \itemch $\overrightarrow{MA}+\overrightarrow{MB}+\overrightarrow{MC}=\overrightarrow0$ khi $M(2;1)$ đúng.
		\end{itemchoice}
	} 
\end{ex}

% nguồn: MaHoaNH_DungSai_Lop10
\begin{ex}%[0H9H1-1]%[0H9H5-2]%[0H9H3-2]%[ma:15719a6c]
	 Trong mặt phẳng tọa độ $Oxy$ cho hai điểm $A(4; 0)$ và $B(0; 2)$.
	 \choiceTF
		{\True Phương trình elip ở dạng chính tắc đi qua $A$ và $B$ là \True $  \dfrac{x^2}{16}+\dfrac{y^2}{4}=1 $}
		{\True Trung điểm của $AB$ là  $M(2;1)$}
		{Phương  trình đường thẳng $AB$ là $x+2y-4=0$}
		{\True$AB=2\sqrt{5}$} 
	 \loigiai{
		\begin{itemize}
			\item Giả sử $\dfrac{x^2}{a^2}+\dfrac{y^2}{b^2}=1$ ($a,b>0$)là phương trình chính tắc  của Elip đi qua $A$ và $B$.\\
			Ta  có
			$$\heva{&\dfrac{4^2}{a^2}+\dfrac{0^2}{b^2}=1\\&\dfrac{0^2}{a^2}+\dfrac{2^2}{b^2}=1}\Leftrightarrow a^2=16,\ b=4.$$ 
			Vậy phương trình chính tắc  của Elip đi qua $A$ VÀ $B$ là $\dfrac{x^2}{16}+\dfrac{y^2}{4}=1 $.
			\item Trung điểm $M$ của $AB$ có tọa  độ là
			 $$\heva{&\dfrac{4+0}{2}=2\\&y_M=\dfrac{0+2}{2}=1.}$$
			Suy ra $M(2;1)$. 
			\item Thay tọa độ của điểm $A(4;0)$ và $B(0;2)$ vào phương trình $x+2y-4=0$ ta thấy tọa độ điểm $A$ và $B$ đều thỏa mãn phương trình nên phương trình $x+2y-4=0$ là phương  trình đường thẳng $AB$.
			 \item Ta có $AB=\sqrt{(0-4)^2+(2-0)}=\sqrt{5}$.
		\end{itemize}
		\begin{itemchoice}
			\itemch  Đúng.
			 \itemch Đúng.
			 \itemch Đúng.
			 \itemch Sai.
		\end{itemchoice}
	} 
\end{ex}


%%%%%%%%%%%%%%%% TRẢ LỜI NGẮN %%%%%%%%%%%%%%%%

% nguồn: MaHoaNH_TraLoiNgan_Lop10
\begin{bt}%[0H9N1-1]%[0D2V2-3]%[ma:86ce3741]
	Một phân xưởng có hai máy đặc chủng $M_1$, $M_2$ sản xuất hai loại sản phẩm kí hiệu là $A$ và $B$. Một tấn sản phẩm loại $A$ lãi $2$ triệu đồng, một tấn sản phẩm loại $B$ lãi $1{,}6$ triệu đồng. Muốn sản xuất một tấn sản phẩm loại $A$ phải dùng máy $M_1$ trong $3$ giờ và máy $M_2$ trong $1$ giờ. Muốn sản xuất một tấn sản phẩm loại $B$ phải dùng máy $M_1$ trong $1$ giờ và  máy $M_2$ trong $1$ giờ. Một máy không thể dùng để sản xuất đồng thời hai loại sản phẩm. Máy $M_1$ làm việc không quá $6$ giờ một ngày, máy $M_2$ làm việc không quá $4$ giờ một ngày. Hỏi số tiền lãi lớn nhất mà phân xưởng này có thể thu được trong một ngày là bao nhiêu triệu đồng?
	 \shortans{6,8}
	 \loigiai{
	\immini{ 
		Gọi $x$ là số tấn sản phẩm loại $A$ và $y$ là số tấn  sản phẩm loại $B$.\\
		Số tiền lãi thu được  là $L=2x+1{,}6y$ triệu đồng.\\
		Theo đề bài  ta có hệ
		$$\heva{&3x+y\le 6\\&x+y\le  4\\&x\ge 0\\&y\ge 0.}$$
		Miền trong của tứ giác $OCDE$, kể cả biên, là biểu diễn miền nghiệm của hệ bất phương trình trên, với $O(0;0)$,  $C(0;4)$, $D(1;3)$, $E(2;0)$.}{
		\begin{tikzpicture}[line cap=round,line join=round,font=\footnotesize,>=stealth,scale=0.7]
				\tikzset{label style/.style={font=\footnotesize}} 
				\draw[->] (-1.25,0)--(5.25,0) node[below  right] {$x$};
				\draw[->]  (0,-1.25)--(0,7.25) node[right] {$y$};
				 \node (0,0) [below left]{$O$};
				\foreach \x  in {-1,1,2,3,4,5}
				\draw[shift={(\x,0)},color=black] (0pt,2pt)--(0pt,-2pt); 
				\node at  (1.75,-0.4) {$E$};
				\node at (.75,-0.4)  {$1$};
				 \node at (4.25,0.4) {$4$};
				\node at  (1.25,3.25) {$D$};
				\foreach \y  in {-1,1,2,3,4,5,6,7}
				\draw[shift={(0,\y)},color=black]  (2pt,0pt)--(-2pt,0pt);
				\node at  (0.5,6.25) {$6$};
				\node at (-0.5,4)  {$C$};
				 \draw [domain=-0.2:2.3,samples=100] plot (\x,{-3*\x+6});
				\node  at (2,4.5) {$3x+y=6$};
				\draw  [domain=-.5:4.5,samples=100] plot (\x,{-\x+4});
				\node  at (4.5,-1) {$x+y=4$};
				 \path
					(0,0) coordinate  (O)
					(0,4) coordinate (C) 
					(2,0)  coordinate (E)
					 (1,3) coordinate (D);
				\fill[pattern=north east lines]  (O)--(C)--(D)--(E)--cycle;
				\draw  (O)--(C)--(D)--(E)--cycle;
			\end{tikzpicture}
		 }
		\noindent Biểu thức $L$ chỉ đạt  giá trị lớn nhất tại một trong bốn đỉnh $O$, $C$, $D$, $E$.\\
		 Tại $O$, ta có $L=0$.\\
		Tại  $C$, ta có $L=6{,}4$.\\
		 Tại $D$, ta có $L=6{,}8$.\\
		Tại $E$, ta có  $L=4$.\\
		Vậy số tiền lớn nhất mà phân xưởng này có thể thu được trong một ngày là $6{,}8$ triệu đồng khi  sản xuất $1$ tấn sản phẩm loại $A$ và $3$ tấn sản phẩm loại $B$.
	} 
\end{bt}

% nguồn: MaHoaNH_TraLoiNgan_Lop10
\begin{bt}%[0H9V1-1]%[ma:a8a6d73a]
	Trong mặt phẳng tọa độ $Oxy$, cho điểm  $M\left(1;-1\right)$ và hai đường thẳng $\Delta\colon x+2y+4=0$ và $d\colon x-y+1=0$. Một đường thẳng qua $M$ cắt $\Delta$ và $d$ lần lượt tại $A$ và $B$ sao cho $M$ là trung điểm của $AB$. Khi đó độ dài $AB$ là bao nhiêu? (kết quả làm tròn đến hàng phần trăm)
	 \shortans{4,72}
	 \loigiai{
		Ta có $A\left(-2a-4;a\right)\in\Delta$,  $B\left(b-1;b\right)\in d$.\\
		Vì $M\left(1;-1\right)$ là trung điểm của  $AB$ nên $\heva{&-2a-4+b-1=2\\ &a+b=-2}\Leftrightarrow\heva{&2a-b=-7\\ &a+b=-2} \Leftrightarrow \heva{&a=-3\\ &b=1}$.\\
		Do dó  $A\left(2;-3\right)$, $B\left(0;1\right)$. Vậy $AB=2\sqrt{5}$.
	 }
\end{bt}


%%%%%%%%%%%%%%%% TỰ LUẬN %%%%%%%%%%%%%%%%

% nguồn: MaHoaNH_TuLuan_Lop10
\begin{bt}%[0H9B1-1]%[ma:35178de4]
	Trong mặt phẳng $Oxy$  cho $A(5;3)$, $B(3;7)$, $C(-2;-8)$. Tìm tọa độ điểm $D$ sao cho $ABCD$ là hình bình hành.
	 \loigiai{
		Gọi $D(x;y)$. Khi đó $\overrightarrow{AD}=(x-5;y-3)$,  $\overrightarrow{BC}=(-5;-15)$.\\
		Tứ giác $ABCD$ là hình bình hành, suy  ra\\
		 $$\overrightarrow{AD}=\overrightarrow{BC}\Leftrightarrow\heva{&x-5=-5\\&y-3=-15}\Leftrightarrow\heva{&x=0\\&y=-12.}$$
		 Vậy $D(0;-12)$.
	} 
\end{bt}
